\documentclass[11pt]{article}

\usepackage[a4paper,margin=1in]{geometry}
\usepackage{amsmath,amssymb,amsthm,mathtools,bm}
\usepackage{microtype}
\usepackage{enumitem}
\usepackage{hyperref}
\usepackage[nameinlink,capitalise]{cleveref}

\hypersetup{colorlinks=true,linkcolor=blue,citecolor=blue,urlcolor=blue}
\setlist{nosep}

\newtheorem{theorem}{Theorem}[section]
\newtheorem{proposition}[theorem]{Proposition}
\newtheorem{lemma}[theorem]{Lemma}
\newtheorem{corollary}[theorem]{Corollary}
\newtheorem{remark}[theorem]{Remark}

\newcommand{\nbar}{\bar n}
\newcommand{\cA}{\mathcal A}
\newcommand{\cO}{\mathcal O}
\newcommand{\cH}{\mathcal H}
\newcommand{\one}{\mathbf 1}
\newcommand{\supp}{\operatorname{supp}}
\newcommand{\diver}{\operatorname{div}}

\title{High-Filling Bose--Hubbard Models on Graphs:\\
Uniform Rotor Comparison, Infrared-Regular Normal Forms,\\
and a Fourth-Order Obstruction}

\author{
Jizheng Chen$^{1,*}$ \qquad Tianrun Yu$^{2,*}$ \qquad Chaoyu Hu$^{3,\dagger}$\\[0.7em]
\small $^1$Department of Mathematics, Chang'an University, Xi'an, China\\
\small $^2$College of Control Science and Engineering, Zhejiang University, Hangzhou, China\\
\small $^3$School of Automation, Hangzhou Dianzi University, Hangzhou, China\\[0.5em]
\small $^*$These authors contributed equally.\\
\small $^\dagger$Corresponding author: \texttt{evaristebernhardwiener@gmail.com}
}
\date{}

\begin{document}
\maketitle

\begin{abstract}
The high-filling Bose--Hubbard model is commonly reduced to a quantum rotor by replacing the bosonic amplitude by its mean value. We establish a finite-graph version of this reduction that keeps the physical one-sided number space and does not impose an a priori cutoff on number fluctuations. The resulting quadratic-form comparison is uniform in graph size and controls number moments, bad-number tails, and the error made when the normalized bosonic annihilation operator is replaced by the phase shift. On a symmetric finite number window we derive the complete second-order normal form. The particle--hole odd first-order term admits an exact completion of squares and satisfies a graph-uniform resolvent estimate, while the direct second-order term is a total double commutator plus a phase-force square. The induced second-order phase potential has an explicit neutral-character expansion and can be compared, through Ginibre monotonicity, with a standard rotor of slightly reduced stiffness. Finally, a three-site classical principal-symbol calculation shows that an all-order ferromagnetic closure fails at fourth order. These results rigorously delimit both the validity and the structural limitations of the high-filling rotor approximation. No thermodynamic Bose--Einstein condensation statement is claimed.
\end{abstract}

\noindent\textbf{Keywords:} Bose--Hubbard model; quantum rotor; high filling; operator comparison; infrared estimate; double commutator; Griffiths--Ginibre inequality.

\section{Introduction}

The Bose--Hubbard Hamiltonian is a standard lattice model for interacting bosons and for the superfluid--Mott transition \cite{Fisher1989}. At a large integer filling $\nbar$, one writes the local occupation as $\nbar+M_x$ and heuristically replaces the normalized annihilation operator by a phase shift. This leads to the compact $U(1)$ quantum rotor. The reduction is widely used in effective-action calculations and in the analysis of Josephson-junction-type systems; see, for example, \cite{PolakKopec2007,Teichmann2009}. Its rigorous content, however, depends on what is being compared: energies, quadratic forms, correlation functions, or low-energy spectral data.

The purpose of this paper is deliberately narrower than a derivation of thermodynamic long-range order. We isolate a package of finite-graph statements that can be proved without assuming condensation and without introducing a number cutoff in the basic comparison. We then compute the finite-window normal form through second order and identify algebraic mechanisms that prevent the most obvious infrared losses. The same calculation also exhibits a sharp limitation: the generalized ferromagnetic structure visible at second order does not persist to all orders.

Our main results are as follows.
\begin{enumerate}
    \item On the physical half-line $M_x\ge -\nbar$, the amplitude defect $A_x-S_x$ is bounded pointwise by $|M_x|/\nbar$. This yields a graph-uniform lower form bound by a quantum rotor with slightly weakened stiffness, together with number-moment and tail estimates.
    \item On a symmetric finite number window, the hopping term has a complete second-order expansion. The first-order particle--hole odd term is an anticommutator with a real phase field; it can be completed into a magnetic square and satisfies a resolvent estimate whose constant is independent of graph volume and of the smallest nonzero graph momentum.
    \item The direct second-order term is a total double commutator plus the square of the phase force. After combining it with the square-completion correction, the remaining phase potential is relatively bounded by the original rotor potential.
    \item On regular graphs, the second-order phase potential has a neutral-character expansion with nonnegative couplings. Standard Ginibre--Griffiths monotonicity therefore compares it with a standard rotor having renormalized stiffness.
    \item A three-site classical elimination shows that at order $\nbar^{-4}$ a character with the opposite sign appears. Thus a naive all-order ferromagnetic expansion is already false at the principal-symbol level.
\end{enumerate}

The work does \emph{not} prove convergence of all correlation functions from the Bose--Hubbard model to the rotor, a volume-uniform Gaussian domination estimate for the full amplitude model, or a new Bose--Einstein condensation regime. These limitations are part of the mathematical statement rather than technical qualifications hidden in the proofs.

\paragraph{Relation to earlier work.}
The high-filling quantum-rotor description of the Bose--Hubbard model is standard in the physics literature \cite{PolakKopec2007,Teichmann2009}. Our contribution is not the heuristic mapping itself, but the finite-graph operator estimates and the explicit protected normal form. The correlation comparison in \cref{sec:characters} uses the general Ginibre framework \cite{Ginibre1970,FrohlichPark1978,Miyao2016} and related quantum-rotor positivity methods; it is quoted as a standard input rather than claimed as a new correlation inequality. The estimates are also conceptually distinct from asymptotic localization results such as \cite{BolsDeRoeck2018}.

\section{Models and notation}

Let $G=(V,E)$ be a finite undirected graph of maximum degree $\Delta$. When regularity is needed, the common degree is denoted by $q$. Each undirected edge is assigned an arbitrary orientation $e=(x,y)$; all final expressions are orientation independent.

The compact rotor at one site acts on $\ell^2(\mathbb Z)$. We use the number operator
\[
M\lvert m\rangle=m\lvert m\rangle
\]
and the bilateral shift
\[
S\lvert m\rangle=\lvert m-1\rangle.
\]
On an edge $e=(x,y)$, set
\begin{equation}
T_e=S_x^\dagger S_y,
\qquad
\cA_e=2-T_e-T_e^\dagger=2-2\cos(\theta_x-\theta_y).
\label{eq:Aedge}
\end{equation}
The rotor Hamiltonian is
\begin{equation}
H_0=H_c+V,
\qquad
H_c=\frac U2\sum_x M_x^2,
\qquad
V=J\sum_e\cA_e.
\label{eq:H0}
\end{equation}

For the physical high-filling block model, the local number deviation satisfies $M_x\ge-\nbar$. The lowering shift is the partial isometry
\[
S\lvert-\nbar\rangle=0,
\qquad
S\lvert m\rangle=\lvert m-1\rangle\quad(m>-\nbar).
\]
Define
\begin{equation}
A_x=S_x\sqrt{1+M_x/\nbar},
\label{eq:Ax}
\end{equation}
and
\begin{equation}
H_{\nbar}=\frac U2\sum_xM_x^2
+J\sum_{e=(x,y)}(A_x-A_y)^\dagger(A_x-A_y).
\label{eq:Hn}
\end{equation}
On a $q$-regular graph and in the zero-total-deviation sector $\sum_xM_x=0$, the onsite part coming from $\sum_e(A_x^\dagger A_x+A_y^\dagger A_y)$ is an additive constant. This is the sector in which the comparison with $H_0$ will be used.

\section{A cutoff-free rotor comparison}

\begin{lemma}[Single-site amplitude defect]
Let $D=A-S$ on the physical one-site Hilbert space. Then
\begin{equation}
D^\dagger D\le \frac{M^2}{\nbar^2}.
\label{eq:defect}
\end{equation}
\end{lemma}

\begin{proof}
Away from the bottom state,
\[
D=S\bigl(\sqrt{1+M/\nbar}-1\bigr).
\]
For $t\ge-1$,
\[
(\sqrt{1+t}-1)^2
=\frac{t^2}{(\sqrt{1+t}+1)^2}
\le t^2.
\]
At the bottom state the left-hand side vanishes, so the same spectral inequality holds everywhere.
\end{proof}

\begin{theorem}[Uniform quadratic-form comparison]
For every $0<\kappa<1$,
\begin{equation}
H_{\nbar}\ge
\left(\frac U2-\frac{2J\Delta}{\kappa\nbar^2}\right)\sum_xM_x^2
+J(1-\kappa)Q_{\mathrm{ph}},
\label{eq:formcomparison}
\end{equation}
where
\begin{equation}
Q_{\mathrm{ph}}:=\sum_{e=(x,y)}(S_x-S_y)^\dagger(S_x-S_y).
\end{equation}
\end{theorem}

\begin{proof}
Write
\[
A_x-A_y=(S_x-S_y)+(D_x-D_y).
\]
The elementary inequality
\[
\|u+v\|^2\ge(1-\kappa)\|u\|^2-\kappa^{-1}\|v\|^2
\]
gives the phase term in \eqref{eq:formcomparison}. Moreover,
\[
\sum_e(D_x-D_y)^\dagger(D_x-D_y)
\le 2\Delta\sum_xD_x^\dagger D_x
\le \frac{2\Delta}{\nbar^2}\sum_xM_x^2
\]
by \eqref{eq:defect}. Substitution into \eqref{eq:Hn} proves the claim.
\end{proof}

A convenient consequence is obtained when $\nbar^2\ge8J\Delta/U$. Choosing
\[
\kappa=\frac{8J\Delta}{U\nbar^2}
\]
yields
\begin{equation}
H_{\nbar}\ge
\frac U4\sum_xM_x^2
+J\left(1-\frac{8J\Delta}{U\nbar^2}\right)Q_{\mathrm{ph}}.
\label{eq:convenientform}
\end{equation}

\begin{corollary}[Moments, tails, and observable replacement]
Assume that $G$ is $q$-regular, work in the sector $\sum_xM_x=0$, and let $\Psi$ be a ground state of $H_{\nbar}$. Under $\nbar^2\ge8J\Delta/U$,
\begin{align}
\frac1{|V|}\sum_x\langle M_x^2\rangle_\Psi
&\le \frac{4\Delta J}{U},
\label{eq:moment}\\
\frac1{|V|}\sum_x\langle(A_x-S_x)^\dagger(A_x-S_x)\rangle_\Psi
&\le \frac{4\Delta J}{U\nbar^2},
\label{eq:obsreplace}\\
\frac1{|V|}\sum_x\langle\one_{\{|M_x|>K\}}\rangle_\Psi
&\le \frac{4\Delta J}{UK^2}.
\label{eq:tails}
\end{align}
\end{corollary}

\begin{proof}
Use the product number state $M_x=0$ as a trial state. Each edge contributes at most $2J$, hence
\[
E_0\le2J|E|\le J\Delta|V|.
\]
Equation \eqref{eq:convenientform} then gives \eqref{eq:moment}. Equation \eqref{eq:obsreplace} follows from \eqref{eq:defect}, and \eqref{eq:tails} follows from
\[
\one_{\{|M|>K\}}\le M^2/K^2.
\]
\end{proof}

The estimates are uniform in $|V|$ and do not use a spectral gap. They are form comparisons; they do not by themselves imply convergence of all low-energy correlations to those of the compact rotor.

\section{Finite-window normal form}

Fix $K\ge1$ and let $P_K$ project onto
\[
\operatorname{span}\{\lvert m\rangle:-K\le m\le K\}.
\]
Write $P_K^{(V)}$ for its tensor product over the graph. All identities in this section concern the compressed matrix elements
\[
P_K^{(V)}H_{\nbar}P_K^{(V)}
\]
of the original physical operators. We do not rebuild the Hamiltonian from $P_KA_xP_K$, because that different operation produces number-window boundary projectors. Assume $\nbar\ge2(K+1)$, so the Taylor expansion is uniform on the window. For $e=(x,y)$, define
\begin{equation}
q_e=M_x+M_y,
\qquad
d_e=M_x-M_y,
\qquad
T_e=S_x^\dagger S_y.
\label{eq:qd}
\end{equation}

\begin{lemma}[Second-order hopping expansion]
On the compressed two-site window,
\begin{equation}
A_x^\dagger A_y
=T_e\left[1+\frac{q_e+1}{2\nbar}-\frac{(d_e+1)^2}{8\nbar^2}\right]+R_e,
\label{eq:hoppingexpansion}
\end{equation}
with
\begin{equation}
\|R_e\|\le C\frac{(K+1)^3}{\nbar^3}.
\label{eq:remainderedge}
\end{equation}
\end{lemma}

\begin{proof}
The nonzero matrix element of $P_K^{(V)}A_x^\dagger A_yP_K^{(V)}$ equals that of the compressed $T_e$ multiplied by
\[
\sqrt{\left(1+\frac{M_x+1}{\nbar}\right)\left(1+\frac{M_y}{\nbar}\right)}.
\]
The claimed expression is the Taylor expansion of $\sqrt{(1+s)(1+t)}$ at $(0,0)$, whose linear part is $(s+t)/2$ and quadratic part is $-(s-t)^2/8$. The third derivatives are uniformly bounded for $|s|,|t|\le1/2$. Both sides vanish on transitions that leave the window. Thus the identity concerns the compressed product $A_x^\dagger A_y$, not a product of separately compressed amplitudes.
\end{proof}

\begin{theorem}[Complete second-order Hamiltonian expansion]
Let
\[
H_{\nbar,K}:=P_K^{(V)}H_{\nbar}P_K^{(V)}.
\]
If $G$ is $q$-regular and we restrict to $\sum_xM_x=0$, then, up to an additive constant,
\begin{equation}
H_{\nbar,K}=H_{0,K}+\nbar^{-1}H_1+\nbar^{-2}H_2+R_3,
\label{eq:normalform}
\end{equation}
where
\begin{align}
H_1&=\frac12K_J+D_{\mathrm{odd}},
&K_J&=-J\sum_e(T_e+T_e^\dagger),
\label{eq:H1}\\
D_{\mathrm{odd}}
&=-\frac J2\sum_eq_e(T_e+T_e^\dagger)
=\frac J2\sum_eq_e\cA_e,
\label{eq:Dodd}\\
H_2&=\frac J8\sum_e\left[T_e(d_e+1)^2+(d_e+1)^2T_e^\dagger\right],
\label{eq:H2}\\
\|R_3\|&\le CJ|E|\frac{(K+1)^3}{\nbar^3}.
\label{eq:R3}
\end{align}
\end{theorem}

\begin{proof}
Compress \eqref{eq:Hn} as a whole and insert \eqref{eq:hoppingexpansion} and its adjoint into the cross-hopping terms. The onsite identity
\[
P_K^{(V)}A_x^\dagger A_xP_K^{(V)}
=P_K^{(V)}\left(1+\frac{M_x}{\nbar}\right)P_K^{(V)}
\]
remains exact; this is why the Hamiltonian is not reconstructed from separately compressed amplitudes. Since $q_e$ commutes with $T_e$ and
\[
\sum_eq_e=q\sum_xM_x=0,
\]
the first-order coefficient separates into $K_J/2$ and $D_{\mathrm{odd}}$. The second-order coefficient is \eqref{eq:H2}; summing the edge remainders gives \eqref{eq:R3}.
\end{proof}

\begin{proposition}[Particle--hole cancellation]
On the symmetric window define $C\lvert m\rangle=\lvert-m\rangle$. Then
\[
CH_0C^*=H_0,
\qquad
CK_JC^*=K_J,
\qquad
CD_{\mathrm{odd}}C^*=-D_{\mathrm{odd}}.
\]
Consequently, for every particle--hole even source $X$ and every $\beta>0$,
\begin{equation}
\langle D_{\mathrm{odd}}\rangle_{\beta,H_0-hX}=0.
\label{eq:phcancellation}
\end{equation}
\end{proposition}

\begin{proof}
One has $CMC^*=-M$ and $CSC^*=S^\dagger$. Thus $T_e$ is exchanged with $T_e^\dagger$, $\cA_e$ is fixed, and $q_e$ changes sign. The Gibbs density of $H_0-hX$ commutes with $C$, so the trace of the odd operator $D_{\mathrm{odd}}$ vanishes.
\end{proof}

\section{The odd term: completion of squares and a graph-uniform resolvent bound}

Define the real multiplication operators
\begin{equation}
b_x=\frac J2\sum_{e\ni x}\cA_e.
\label{eq:bx}
\end{equation}

\begin{proposition}[Phase-defect representation]
In the $q$-regular zero-charge sector,
\begin{equation}
D_{\mathrm{odd}}=\frac12\sum_x\{b_x,M_x\}.
\label{eq:anticomm}
\end{equation}
\end{proposition}

\begin{proof}
For an edge $e=(x,y)$, the total edge charge $q_e=M_x+M_y$ commutes with $\cA_e$. Summing the two incidence contributions gives
\[
\frac J4\{\cA_e,q_e\}=\frac J2q_e\cA_e,
\]
which is exactly the edge term in \eqref{eq:Dodd}.
\end{proof}

\begin{proposition}[Exact completion of squares]
With $\varepsilon=\nbar^{-1}$,
\begin{equation}
\frac U2\sum_xM_x^2+\varepsilon D_{\mathrm{odd}}
=\frac U2\sum_x\left(M_x+\frac{\varepsilon b_x}{U}\right)^2
-\frac{\varepsilon^2}{2U}\sum_xb_x^2.
\label{eq:squarecompletion}
\end{equation}
\end{proposition}

\begin{proof}
Expand each noncommutative square. The cross term is $\frac\varepsilon2\{M_x,b_x\}$; no commutativity between $M_x$ and $b_x$ is required.
\end{proof}

In the angle representation, \eqref{eq:squarecompletion} is a Schr\"odinger operator with a real magnetic vector potential. The standard diamagnetic inequality therefore compares its ground-state energy with the same scalar potential and zero vector potential. The statement remains valid in the zero-total-charge sector because taking the modulus preserves invariance under common angle shifts.

\begin{theorem}[Composite infrared estimate]
Let $\Omega>0$ be the normalized ground state of $H_0$ in the zero-total-charge sector, set $K_0=H_0-E_0$, and let $K_{0,\perp}^{-1}$ denote the inverse on $\Omega^\perp$. For real edge coefficients $u_e$, define
\begin{equation}
\cO(u)=\sum_{e=(x,y)}u_e(M_x+M_y)\cA_e.
\label{eq:Ou}
\end{equation}
Then
\begin{equation}
\left\langle \cO(u)\Omega,
K_{0,\perp}^{-1}\cO(u)\Omega\right\rangle
\le \frac{4\Delta}{U}\sum_e|u_e|^2
\langle\Omega,\cA_e\Omega\rangle
\le \frac{16\Delta}{U}\sum_e|u_e|^2.
\label{eq:infrared}
\end{equation}
\end{theorem}

\begin{proof}
Set
\[
B_x(\theta)=\sum_{e\ni x}u_e\cA_e(\theta).
\]
Since each $\cA_e$ depends only on $\theta_x-\theta_y$, the full divergence satisfies
\[
\sum_x\partial_{\theta_x}B_x=0.
\]
Hence $\cO(u)=-iB\cdot\nabla$. The ground-state transform gives
\[
\langle f\Omega,K_0f\Omega\rangle
=\frac U2\int|\nabla f|^2\Omega^2.
\]
Integration by parts yields
\[
\langle f\Omega,\cO(u)\Omega\rangle
=\frac i2\int\Omega^2B\cdot\nabla\overline f.
\]
Optimizing the resolvent variational formula gives an upper bound by $(2U)^{-1}\int|B|^2\Omega^2$. Finally,
\[
|B|^2\le2\Delta\sum_e|u_e|^2\cA_e^2,
\qquad
\cA_e^2\le4\cA_e\le16,
\]
which proves \eqref{eq:infrared}.
\end{proof}

The important feature of \eqref{eq:infrared} is the absence of the graph spectral gap. The composite density--phase operator has a built-in transport cancellation; estimating it only by operator norm would lose this structure.

\section{The direct second-order term}

\begin{lemma}[Edge double commutator]
For $C_e=M_x^2+M_y^2$,
\begin{equation}
T_e(d_e+1)^2+(d_e+1)^2T_e^\dagger
=\frac14[C_e,[C_e,T_e+T_e^\dagger]].
\label{eq:doublecommedge}
\end{equation}
\end{lemma}

\begin{proof}
The shift $T_e$ changes $d_e$ by two and changes $C_e$ by $2(d_e+1)$. Thus
\[
[C_e,T_e]=2T_e(d_e+1),
\]
and a second commutator gives $4T_e(d_e+1)^2$. Taking the adjoint gives the $T_e^\dagger$ term with the displayed ordering.
\end{proof}

\begin{theorem}[Global normal form]
The operator $H_2$ in \eqref{eq:H2} satisfies
\begin{equation}
H_2=-\frac1{8U^2}[H_c,[H_c,V]],
\label{eq:H2comm1}
\end{equation}
and
\begin{equation}
H_2=-\frac1{8U^2}[H_0,[H_0,V]]
+\frac1{8U}|\nabla_\theta V|^2.
\label{eq:H2comm2}
\end{equation}
\end{theorem}

\begin{proof}
Equation \eqref{eq:H2comm1} follows by summing \eqref{eq:doublecommedge} and using
\[
V=2J|E|-J\sum_e(T_e+T_e^\dagger).
\]
Since $[H_c,V]=[H_0,V]$,
\[
[H_c,[H_c,V]]=[H_0,[H_0,V]]-[V,[H_0,V]].
\]
For $H_c=-(U/2)\Delta_\theta$,
\[
[V,[H_c,V]]=U|\nabla_\theta V|^2,
\]
which proves \eqref{eq:H2comm2}.
\end{proof}

Combining the phase-force term in \eqref{eq:H2comm2} with the negative scalar correction in \eqref{eq:squarecompletion} defines
\begin{equation}
W_{\mathrm{sc}}
=\frac1{8U}|\nabla_\theta V|^2
-\frac1{2U}\sum_xb_x^2.
\label{eq:Wscdef}
\end{equation}

\begin{proposition}[Relative phase bound]
Pointwise on the phase torus,
\begin{equation}
|W_{\mathrm{sc}}|\le \frac{J\Delta}{U}V.
\label{eq:relativephase}
\end{equation}
\end{proposition}

\begin{proof}
Let
\[
Q_x=\sum_{y\sim x}\left(e^{i(\theta_x-\theta_y)}-1\right).
\]
Then
\[
\partial_{\theta_x}V=2J\,\operatorname{Im}Q_x,
\qquad
b_x=-J\,\operatorname{Re}Q_x,
\]
so
\[
W_{\mathrm{sc}}
=\frac{J^2}{2U}\sum_x
\left[(\operatorname{Im}Q_x)^2-(\operatorname{Re}Q_x)^2\right].
\]
By Cauchy--Schwarz,
\[
|Q_x|^2
\le \deg(x)\sum_{y\sim x}|e^{i(\theta_x-\theta_y)}-1|^2
=\deg(x)\sum_{y\sim x}\cA_{xy}.
\]
Summing over $x$ yields
\[
\sum_x|Q_x|^2\le \frac{2\Delta}{J}V,
\]
which gives \eqref{eq:relativephase}.
\end{proof}

Thus $\nbar^{-2}W_{\mathrm{sc}}$ is relatively of size $J/(U\nbar^2)$ compared with the principal phase potential. The statement is pointwise and volume independent.

\section{Neutral-character structure at second order}
\label{sec:characters}

\begin{proposition}[Character expansion]
If $G$ is $q$-regular, then
\begin{equation}
W_{\mathrm{sc}}
=-\frac{J^2}{4U}\sum_x(Q_x^2+\overline{Q}_x^{\,2}),
\label{eq:WscQ}
\end{equation}
and
\begin{equation}
W_{\mathrm{sc}}
=\frac{J^2q}{U}\sum_x\sum_{y\sim x}\cos(\theta_x-\theta_y)
-\frac{J^2}{2U}\sum_x\sum_{y,z\sim x}\cos(2\theta_x-\theta_y-\theta_z)
+\mathrm{const}.
\label{eq:Wsccharacters}
\end{equation}
\end{proposition}

\begin{proof}
The first formula follows from
\[
(\operatorname{Im}Q)^2-(\operatorname{Re}Q)^2
=-\frac12(Q^2+\overline Q^{\,2}).
\]
Expanding $Q_x^2$ and taking real parts gives
\[
Q_x^2+\overline Q_x^{\,2}
=2\sum_{y,z\sim x}\cos(2\theta_x-\theta_y-\theta_z)
-4q\sum_{y\sim x}\cos(\theta_x-\theta_y)
+2q^2.
\]
\end{proof}

Consequently, up to an additive constant, the phase Hamiltonian through order $\nbar^{-2}$ can be written as
\begin{equation}
V_{\mathrm{eff}}
=J_*\sum_e\cA_e
-\frac{J^2}{2U\nbar^2}
\sum_x\sum_{y,z\sim x}\cos(2\theta_x-\theta_y-\theta_z),
\label{eq:Veff}
\end{equation}
where
\begin{equation}
J_*=J-\frac{qJ^2}{U\nbar^2}.
\label{eq:Jstar}
\end{equation}
Every new character in \eqref{eq:Veff} is neutral: its integer coefficients sum to zero. Its Hamiltonian coefficient has the ferromagnetic sign.

\begin{proposition}[Standard Ginibre monotonicity for neutral characters]
Consider a finite-volume quantum rotor
\begin{equation}
H(\bm J)
=\frac U2\sum_xM_x^2
-\sum_\gamma J_\gamma\cos(\gamma\cdot\theta),
\qquad
\sum_x\gamma_x=0,
\qquad
J_\gamma\ge0.
\label{eq:ginibresetup}
\end{equation}
For neutral $\eta$, the finite-temperature expectation
$\langle\cos(\eta\cdot\theta)\rangle$ is nondecreasing in every coupling $J_\gamma$. The same is true for zero-temperature limit states.
\end{proposition}

\begin{proof}[Standard input]
A Trotter discretization in imaginary time has nonnegative Fourier coefficients both for the rotor heat kernel and for $\exp(tJ_\gamma\cos(\gamma\cdot\theta))$. Ginibre's positive-function cone gives the second correlation inequality for the resulting classical compact-abelian spin system \cite{Ginibre1970,FrohlichPark1978}. Passing to the Trotter and zero-temperature limits gives the statement; see also quantum formulations such as \cite{Miyao2016,Miyao2015,BjornbergUeltschi2023}. We use this as an external positivity theorem rather than as a new result.
\end{proof}

\begin{corollary}[Comparison with a weakened standard rotor]
If $J_*>0$, then for every pair $x,y$,
\begin{equation}
\langle\cos(\theta_x-\theta_y)\rangle_{V_{\mathrm{eff}}}
\ge
\langle\cos(\theta_x-\theta_y)\rangle_{\text{standard rotor of stiffness }J_*}.
\label{eq:rotorcomparison}
\end{equation}
\end{corollary}

\begin{proof}
Start from the standard rotor with stiffness $J_*$ and turn on the nonnegative neutral couplings in \eqref{eq:Veff} one at a time. The preceding monotonicity statement yields \eqref{eq:rotorcomparison}.
\end{proof}

\section{A fourth-order obstruction to all-order ferromagnetic closure}

The preceding section concerns the second-order effective phase potential. It is natural, but incorrect, to conjecture that all higher-order density eliminations preserve the same sign structure. The following three-site principal-symbol calculation gives a strict counterexample.

Consider the path $0$--$1$--$2$. Set $\theta_0=0$, $\theta_1=a$, $\theta_2=b$, $c_1=\cos a$, and $c_2=\cos(b-a)$. Under $m_0+m_1+m_2=0$, define
\begin{equation}
\cH_\varepsilon
=\frac U2\sum_{i=0}^2m_i^2
-K_0\left[
\sqrt{(1+\varepsilon m_0)(1+\varepsilon m_1)}\,c_1
+\sqrt{(1+\varepsilon m_1)(1+\varepsilon m_2)}\,c_2
\right].
\label{eq:classical3site}
\end{equation}

\begin{theorem}[Three-site fourth-order obstruction]
The local minimizer near $m=0$ has
\begin{align}
m_0&=\frac{K_0(c_1-2c_2)}{6U}\varepsilon
+\frac{K_0^2c_1c_2}{8U^2}\varepsilon^3
+O(\varepsilon^5),
\label{eq:m0exp}\\
m_1&=\frac{K_0(c_1+c_2)}{6U}\varepsilon
-\frac{K_0^2c_1c_2}{4U^2}\varepsilon^3
+O(\varepsilon^5),
\label{eq:m1exp}
\end{align}
and the minimized energy equals
\begin{align}
E_{\mathrm{eff}}
={}&-K_0(c_1+c_2)
+\frac{K_0^2\varepsilon^2}{12U}
(-c_1^2+c_1c_2-c_2^2)
\notag\\
&+\frac{K_0^3\varepsilon^4}{32U^2}
c_1c_2(c_1+c_2)
+O(\varepsilon^6).
\label{eq:Eeff4}
\end{align}
The $\varepsilon^4$ term contains
\begin{equation}
+\frac{K_0^3\varepsilon^4}{128U^2}\cos(a+b),
\label{eq:wrongsign}
\end{equation}
which has the opposite sign from a ferromagnetic Hamiltonian term $-\lambda\cos(a+b)$ with $\lambda\ge0$.
\end{theorem}

\begin{proof}
Write $m_2=-m_0-m_1$ and solve the two stationary equations with the ansatz
\[
m_0=a_1\varepsilon+a_3\varepsilon^3+O(\varepsilon^5),
\qquad
m_1=b_1\varepsilon+b_3\varepsilon^3+O(\varepsilon^5).
\]
Coefficient matching gives \eqref{eq:m0exp}--\eqref{eq:m1exp}. Substitution yields \eqref{eq:Eeff4}. Finally,
\begin{align*}
c_1c_2(c_1+c_2)
={}&\frac12\cos a+\frac12\cos(a-b)
+\frac14\cos(a+b)+\frac14\cos(a-2b)\\
&+\frac14\cos(3a-b)+\frac14\cos(3a-2b),
\end{align*}
which proves \eqref{eq:wrongsign}.
\end{proof}

The theorem is intentionally stated as a classical principal-symbol obstruction. It does not by itself identify the complete fourth-order quantum effective Hamiltonian. It does rule out arguments asserting that the density-amplitude expansion is sign-preserving at every order for purely algebraic reasons.

\begin{lemma}[Neutral characters are controlled by local gradients]
Let $\gamma\in\mathbb Z^V$ satisfy $\sum_x\gamma_x=0$ and suppose $\supp\gamma$ lies in a connected finite subgraph $X$. Then
\begin{equation}
1-\cos(\gamma\cdot\theta)
\le C_{\gamma,X}\sum_{e\subset X}
\bigl[1-\cos(\nabla_e\theta)\bigr].
\label{eq:neutralgradient}
\end{equation}
\end{lemma}

\begin{proof}
Choose an integer flow $c_e$ on a spanning tree of $X$ with $\diver c=\gamma$. Then
\[
\gamma\cdot\theta=\sum_ec_e\nabla_e\theta.
\]
The inequalities
\[
|e^{in\phi}-1|\le |n|\,|e^{i\phi}-1|
\]
and a telescoping product bound give
\[
|e^{i\gamma\cdot\theta}-1|
\le \sum_e|c_e|\,|e^{i\nabla_e\theta}-1|.
\]
Squaring and applying Cauchy--Schwarz proves \eqref{eq:neutralgradient}.
\end{proof}

\section{Scope, limitations, and possible extensions}

The results establish a rigorous and graph-uniform structural approximation, but they stop short of a complete low-energy equivalence. Three distinctions are essential.
\begin{itemize}
    \item The cutoff-free comparison controls quadratic forms, moments, and tails. It does not imply convergence of the full spectrum or of arbitrary correlation functions.
    \item The second-order character comparison applies to the effective compact-rotor Hamiltonian. It is not a Gaussian domination theorem for the original amplitude model.
    \item The fourth-order obstruction concerns the classical principal symbol. It invalidates an all-order sign conjecture but does not preclude perturbative control by other mechanisms.
\end{itemize}

A stronger theorem would require one of two additional ingredients: either a volume-uniform source estimate that transfers rotor infrared bounds to the amplitude model, or a direct norm-resolvent/low-energy-subspace comparison in a parameter regime with controlled number tails. Neither statement follows from the present normal form alone.

The package nevertheless supplies a reusable diagnostic for other gapless bosonic lattice models: odd density--phase couplings should first be tested for an anticommutator representation and a ground-state-transform estimate, while direct even corrections should be tested for total commutators before they are bounded in operator norm.

\appendix

\section{Coefficient audit}

For convenience, we collect the normalization used throughout the paper:
\begin{align}
H_{\nbar}
&=\frac U2\sum_xM_x^2
+J\sum_e(A_x-A_y)^\dagger(A_x-A_y),
\label{eq:audit1}\\
H_0
&=\frac U2\sum_xM_x^2+J\sum_e\cA_e,
\qquad
\cA_e=2-T_e-T_e^\dagger,
\label{eq:audit2}\\
D_{\mathrm{odd}}
&=\frac J2\sum_eq_e\cA_e,
\qquad
b_x=\frac J2\sum_{e\ni x}\cA_e,
\label{eq:audit3}\\
H_2
&=\frac J8\sum_e\left[T_e(d_e+1)^2+(d_e+1)^2T_e^\dagger\right],
\label{eq:audit4}\\
H_2
&=-\frac1{8U^2}[H_0,[H_0,V]]
+\frac1{8U}|\nabla_\theta V|^2,
\label{eq:audit5}\\
W_{\mathrm{sc}}
&=-\frac{J^2}{4U}\sum_x(Q_x^2+\overline Q_x^{\,2}),
\qquad
J_*=J-\frac{qJ^2}{U\nbar^2}.
\label{eq:audit6}
\end{align}
Changing the coefficient of the phase-gradient term from $J$ to $J/2$ requires a global replacement of $J$ in every hopping, force-square, and stiffness-renormalization formula. Mixing the two conventions produces the factor-of-two inconsistencies that motivated the normalization audit underlying this version.

\section{Statements not claimed}

For clarity, this paper does not claim any of the following:
\begin{itemize}
    \item a new thermodynamic Bose--Einstein condensation exponent or density regime;
    \item volume-uniform Gaussian domination for the original high-filling amplitude model;
    \item removal of the finite number window in the second-order operator-norm expansion;
    \item an all-order ferromagnetic effective Hamiltonian;
    \item identification of the complete fourth-order quantum normal form.
\end{itemize}

\bibliographystyle{plain}
\bibliography{references}

\end{document}
